% Preamble of observational_classification.tex: packages, colours, the TikZ vocabulary of the figures, boxes, macros.
\usepackage[a4paper,margin=2.35cm]{geometry}
\usepackage[T1]{fontenc}
\usepackage[utf8]{inputenc}
\usepackage{lmodern}
\usepackage{microtype}
\usepackage{amsmath,amssymb,amsthm,mathtools}
\usepackage{booktabs,array,multirow}
\usepackage{enumitem}
\usepackage{float}
\usepackage[font=small,labelfont=bf]{caption}
\usepackage{xcolor}
\usepackage{tikz}
\usetikzlibrary{arrows.meta,positioning,calc,backgrounds,fit,shapes.geometric,decorations.pathreplacing}
\usepackage[most]{tcolorbox}
\usepackage[colorlinks=true,linkcolor=oiblue!80!black,citecolor=oigreen!70!black,urlcolor=oiblue!80!black]{hyperref}
\usepackage[capitalise,noabbrev]{cleveref}

% ---------------------------------------------------------------------------------------------------------------
% Colours: the Okabe-Ito palette, which stays distinguishable under the common colour-vision deficiencies.
% ---------------------------------------------------------------------------------------------------------------
\definecolor{oiorange}{RGB}{230,159,0}
\definecolor{oisky}{RGB}{86,180,233}
\definecolor{oigreen}{RGB}{0,158,115}
\definecolor{oiyellow}{RGB}{240,228,66}
\definecolor{oiblue}{RGB}{0,114,178}
\definecolor{oired}{RGB}{213,94,0}
\definecolor{oipurple}{RGB}{204,121,167}
\colorlet{qcol}{oipurple!82!black}      % quantum sources and what they send
\colorlet{ccol}{oiblue}                 % classical data handed out by a source
\colorlet{Cfill}{oisky!55}              % the part C that is split off
\colorlet{Dfill}{oigreen!38}            % the part D that sees C
\colorlet{splitcol}{oiorange}           % the source that is being split
\colorlet{dropcol}{oired}               % the edge that is dropped
\colorlet{premcol}{oiyellow!70}         % premeasurable nodes
\colorlet{neutral}{black!55}            % boxes and discs that carry no status

% ---------------------------------------------------------------------------------------------------------------
% TikZ vocabulary. A structure is drawn with \vn (visible node), \qs (quantum source), \feed (source to members)
% and \de (visible edge).
% ---------------------------------------------------------------------------------------------------------------
\tikzset{
  mdag/.style={>={Stealth[length=2.0mm,width=1.6mm]}, line cap=round, line join=round, font=\small},
  vis/.style={circle, draw=black, line width=0.7pt, fill=white, minimum size=6.6mm, inner sep=0pt, font=\small},
  Cn/.style={vis, fill=Cfill},
  Dn/.style={vis, fill=Dfill},
  pre/.style={vis, fill=premcol},
  src/.style={diamond, aspect=1, draw=qcol!75!black, line width=0.5pt, fill=qcol, minimum size=4.6mm, inner sep=0pt},
  splitsrc/.style={src, draw=splitcol!70!black, fill=splitcol},
  csrc/.style={src, draw=ccol!75!black, fill=ccol!65},
  edge/.style={->, line width=0.8pt, black},
  qedge/.style={->, line width=0.8pt, qcol},
  splitedge/.style={->, line width=0.8pt, splitcol!85!black},
  cedge/.style={->, line width=0.8pt, ccol, densely dashed},
  dropedge/.style={->, line width=0.9pt, dropcol, densely dashed},
  newedge/.style={->, line width=1.0pt, oigreen!75!black},
  faint/.style={opacity=0.28},
  lab/.style={font=\footnotesize, inner sep=1.5pt},
  rel/.style={font=\Large},
  tag/.style={font=\scriptsize\itshape, text=black!70},
}
\newcommand{\vn}[4][vis]{\node[#1] (#2) at (#3) {#4};}
\newcommand{\qs}[3][src]{\node[#1] (#2) at (#3) {};}
\newcommand{\feed}[3][qedge]{\foreach \target in {#3} {\draw[#1] (#2) -- (\target);}}
\newcommand{\de}[3][edge]{\draw[#1] (#2) -- (#3);}
% relation signs between two pictures
\newcommand{\equivsign}{\ensuremath{\boldsymbol{\equiv}}}
\newcommand{\nequivsign}{\ensuremath{\boldsymbol{\not\equiv}}}
\newcommand{\openequivsign}{\ensuremath{\overset{?}{\equiv}}}

% ---------------------------------------------------------------------------------------------------------------
% Boxes
% ---------------------------------------------------------------------------------------------------------------
\tcbset{boxbase/.style={enhanced, breakable, boxrule=0.8pt, arc=1.6mm, left=2.2mm, right=2.2mm, top=1.6mm,
                        bottom=1.6mm, fonttitle=\bfseries\small, before skip=9pt, after skip=9pt}}
\newtcolorbox{rulebox}[1]{boxbase, colback=oiblue!4, colframe=oiblue, title={#1}}
\newtcolorbox{validbox}[1]{boxbase, colback=oigreen!5, colframe=oigreen!85!black, title={\checkmark\enspace #1}}
\newtcolorbox{warnbox}[1]{boxbase, colback=oired!5, colframe=oired, title={$\boldsymbol{\times}$\enspace Not like this: #1}}
\newtcolorbox{openq}[1]{boxbase, colback=oiorange!7, colframe=oiorange!90!black, title={?\enspace Open: #1}}
\newtcolorbox{ideabox}[1]{boxbase, colback=black!3, colframe=black!55, title={#1}}

% ---------------------------------------------------------------------------------------------------------------
% Theorems
% ---------------------------------------------------------------------------------------------------------------
\theoremstyle{plain}
\newtheorem{theorem}{Theorem}
\newtheorem{lemma}[theorem]{Lemma}
\newtheorem{proposition}[theorem]{Proposition}
\newtheorem{corollary}[theorem]{Corollary}
\theoremstyle{definition}
\newtheorem{definition}[theorem]{Definition}
\newtheorem{example}[theorem]{Example}
\theoremstyle{remark}
\newtheorem{remark}[theorem]{Remark}

% ---------------------------------------------------------------------------------------------------------------
% Macros
% ---------------------------------------------------------------------------------------------------------------
\newcommand{\Q}{\mathcal{Q}}                         % quantum model set
\newcommand{\C}{\mathcal{C}}                         % classical model set
\newcommand{\G}{G}
\newcommand{\Pa}{\mathrm{Pa}}
\newcommand{\pa}{\mathrm{pa}}
\newcommand{\Tr}{\operatorname{Tr}}
\newcommand{\Hil}{\mathcal{H}}
\newcommand{\id}{\mathbf{1}}
\newcommand{\ket}[1]{\lvert #1\rangle}
\newcommand{\bra}[1]{\langle #1\rvert}
\newcommand{\proj}[1]{\lvert #1\rangle\!\langle #1\rvert}
\newcommand{\code}[1]{\texttt{#1}}
\newcommand{\str}[1]{\textsf{\small #1}}             % a structure in the repository's notation
\newcommand{\qeq}{\equiv_{\Q}}
\newcommand{\ceq}{\equiv_{\C}}
\newcommand{\AWS}{AWS}
\setlist{itemsep=2pt, topsep=4pt}
\setlength{\parskip}{4pt plus 1pt}
\setlength{\parindent}{0pt}
